\documentclass{article}
\usepackage{amsmath, amssymb, amsthm}
\usepackage{hyperref}
\usepackage{geometry}
\geometry{margin=1in}

\title{Erd\H{o}s Problem \#575}
\date{Last updated: 2025-08-31}
\author{Source: erdosproblems.com}

\begin{document}
\maketitle

\noindent\textbf{Status:} OPEN \quad \textbf{No prize}

\noindent\textbf{Tags:} graph theory, turan number

\medskip

\noindent\textbf{Problem Statement:}

If $\mathcal{F}$ is a finite set of finite graphs then $\mathrm{ex}(n;\mathcal{F})$ is the maximum number of edges a graph on $n$ vertices can have without containing any subgraphs from $\mathcal{F}$. Note that it is trivial that $\mathrm{ex}(n;\mathcal{F})\leq \mathrm{ex}(n;G)$ for every $G\in\mathcal{F}$. Is it true that, for every $\mathcal{F}$, if there is a bipartite graph in $\mathcal{F}$ then there exists some bipartite $G\in\mathcal{F}$ such that\[\mathrm{ex}(n;G)\ll_{\mathcal{F}}\mathrm{ex}(n;\mathcal{F})?\]




A problem of Erd\H{o}s and Simonovits. See also [180] . This problem is #51 in Extremal Graph Theory in the graphs problem collection.


\bigskip
\noindent\small{Source: \url{https://www.erdosproblems.com/575}}
\end{document}
